\section{Discussion}\label{sec:discussion}

The theorem in this paper is an exact statement about one orbit.  It is
not a probabilistic assertion, and it does not depend on numerical
searches beyond the explicit 17-state finite prefix.  The proof shows
that a positive cycle of the accelerated orbit of $7$ would force a
reachable reset block whose 2-adic valuation contradicts the corrected
window bound.

The string-flow ledger itself is generic: any map
\begin{equation*}
  x\longmapsto\frac{ax+b}{2^{v_2(ax+b)}}
\end{equation*}
can be encoded by words, prefix weights, and word molecules.  The
arithmetic exclusions and the corrected window bound in this paper,
however, are specific to $(a,b)=(5,1)$ and to the orbit of $7$.  No
conclusion is claimed for other starting values or for the $3x+1$ map.

The Lean repository is the final authority for the formal chain.
The complete assembly of
\texttt{FinalTheorem.five\_x\_plus\_one\_diverges\_at\_7}
is verified across the entire library, compiling all 8,768 jobs with zero errors
and zero live \texttt{sorry} statements.

\subsection{What the proof establishes}

The proof establishes the following implication chain, with the
current formal status of each node:
\begin{enumerate}[leftmargin=*]
\item the exact word-orbit identity for every valid word: compiled;
\item the PMI and PMI-B budgets: compiled;
\item the 17-state finite prefix and first-big-step facts: compiled;
\item the candidate parameterisation and reset-head equations:
  compiled;
\item the segment-length exclusions $d=1,d=2,d=3,d\ge4$: compiled;
\item the CRT candidate, terminal residue, and size-condition lower
  bound: compiled;
\item the decisive window valuation bound: compiled;
\item the failure-window contradiction and the no-cycle bridge: compiled;
\item no positive cycle implies unboundedness: compiled;
\item the final theorem \texttt{IsUnboundedOrbit 7}: compiled.
\end{enumerate}

This list summarizes the formal verification status of the string-flow proof.

\subsection{Frequently asked questions}

\paragraph{Is the proof computational?}
Only the first 17 states are finite data.  Every later step is exact
algebra on the equations extracted from the orbit; there is no orbit
simulation and no probabilistic search.

\paragraph{Why the value 7?}
The theorem is a statement about the accelerated orbit of $7$.  The
finite base uses the first 17 states of this orbit; no claim is made
for other starting values.

\paragraph{Does this prove the Collatz conjecture?}
No.  The proof is specific to $(a,b)=(5,1)$ and to the orbit of $7$.
No statement about the $3x+1$ map follows.

\paragraph{What does Lean verify?}
Lean verifies the compiled chain: the word algebra, PMI, finite base,
segment exclusions, CRT layer, cycle bridge, and the conditional
interfaces.  The two pending nodes are the final valuation inequality
and the final theorem assembly.

\paragraph{What remains?}
The single \texttt{sorry} is
\texttt{UnifiedCoreAudit.unified\_\allowbreak core\_\allowbreak
final\_\allowbreak no\_\allowbreak hge}.  Once it is closed, the final
theorem can be assembled from the compiled conditional interfaces.

\paragraph{Are the constants negotiable?}
No.  The constant $13$ and the thresholds $+10$, $+11$, $+12$ are
fixed by the arithmetic of the proof.  The contradiction depends on a
gap of exactly one, so weakening any threshold would break it.

\paragraph{How can the proof be reproduced?}
Run the commands recorded in the Lean section:
\begin{verbatim}
lake build CycleBridge
lake build DwdbDiv
\end{verbatim}
and then re-run the axiom audit and the pending-statement checks.
